\section{Conclusion}
\label{sec:conclusion}

A lunar PNT realization fixes a spatial frame and a timescale jointly, and this paper
classifies what such a realization fails to determine about the two together. Over a
nine-parameter coupled datum---translation, scale, rotation, a timescale offset $\dtau$, and
a timescale rate $\dalpha$---we prove a \emph{basis-invariant} classification of the coupled
Fisher null space into spatial, temporal, and genuinely coupled parts, expressed through the
intrinsic projector $\Pnull=\Umat\Umat^\top$ and its off-diagonal block $\Pst$. Under the SI
second and the $c$-defined metre---adopted up front to quarantine the definitional
unit-of-time/scale coupling and keep scale observable---the exact datum gauge of a fully
internal \emph{free-network} one-way lunar problem (mutual common-mode clocks, no external
tie) is \emph{eight}: three translations, three rotations, the offset, and the rate;
anchoring to known orbiter positions and a time reference reduces this to defect zero, the
validated case. Scale is observable, so we never treat it as a free datum direction. A
coupled defect is certified only by $\Pst\neq\bm0$; a non-block-diagonal Fisher is necessary
but not sufficient. External ties lift the temporal defects---a time tie removes $\dtau$, and
the relativistic rate tie removes $\dalpha$ given an independent scale tie---and, once the
rate tie is admitted, the joint gauge acquires a spatial--temporal direction and is
\emph{not} the direct sum of the spatial and temporal gauges.

We quantify the two frame--time couplings and state one of them as an explicit correction.
Coupling~1 (range$\leftrightarrow$clock-offset, unit tie
$1\,\text{m}\leftrightarrow3.336\,\text{ns}=1/c$) is operationally strong but a near-defect,
and at the datum level, in a geometrically diverse one-way network, two-way ranging provides
\emph{no} selective scale$\leftrightarrow$offset decoupling beyond data volume (lift ratio
$\approx1.00$, computed on the modelled network)---the strong degeneracy is a single-user,
poor-geometry effect that does not aggregate. Coupling~2 (relativistic rate$\leftrightarrow$frame) is fundamentally weak: the
frame estimation datum perturbs the lunar timescale rate by $\lesssim\!\sim\!2\times10^{-17}$,
six orders of magnitude below the $\pm1.7\times10^{-11}$ reference-surface definitional band;
the rate is gauge-robust to estimation and gauge-sensitive to definition.

Exactly one claim is externally \Validated{}: the coupled-Fisher linear-algebra pipeline
reproduces an independent numpy/LAPACK computation, to relative and absolute error below
$10^{-3}$, on rows carrying real DE440 lunar-orientation partials, for a full-rank
(defect-$0$) and an exact-rate-defect network. The geometry is \Modelled{}, the coupling
magnitudes are \Modelled{}, and the nonzero-coupling branch of the classifier is \Proven{} by
construction and unit-tested, not externally validated---stated precisely.

Two design laws follow: geometric diversity separates the frame scale from the timescale
offset; and the rate's residual gauge is definitional, not estimational. Together they give
the compact requirement (Theorem~\ref{thm:system}): \emph{a lunar reference-system standard
must fix both a spatial datum and a temporal convention}---neither geometry, the relativistic
model, nor any volume of data can supply what the other must. The proven and cited content
carries over unchanged to any real deployment; before any absolute rate or frame number is
trusted, a real lunar realization must supply flight data and a fixed convention. Every number
here is reproducible from the open engine by a single command.
